% ch01 — Special Relativity: A Working Primer
\chapter{Special Relativity: A Working Primer}
\label{ch:special-relativity}

Everything in this book---the geodesic integrator, the metric functions, the
camera model, the accretion disc renderer---rests on a handful of structures
from special relativity: four-vectors, the Minkowski metric, and the causal
distinction between timelike, null, and spacelike directions.  This
chapter builds those structures from first principles, establishing the index
notation and conventions that we'll use without further ceremony for the
remaining seven hundred-odd pages.
\histnote{Einstein's 1905 paper derived the Lorentz transformations from two
postulates.  We take the geometric view instead: define the metric, and let the
symmetries follow.}

The emphasis is deliberately computational.  We won't trace the historical
path from Michelson--Morley through Einstein's postulates to the Lorentz
transformations---excellent accounts exist elsewhere
\cite{Carroll2004,MTW1973}.  Instead, we'll work in geometrised units ($c = G
= 1$) from the start and focus on the tensorial machinery that appears
directly in the codebase: the metric as a matrix, index operations as linear
algebra, and null vectors as the starting data for ray tracing.

Readers comfortable with special relativity may wish to skim this chapter for
notation and proceed to \cref{ch:differential-geometry}, where we generalise
to curved spacetime.  But the tetrad construction in \cref{sec:sr-lorentz} and
the null-vector parametrisation in \cref{sec:sr-causal} reappear verbatim in
the camera model of \cref{ch:camera-rendering}, so they're worth a careful
look even for the experienced reader.

% ---- Drafted sections ----
% sec:sr-minkowski — The Minkowski Metric, Four-Vectors, and Index Gymnastics
\section{The Minkowski Metric, Four-Vectors, and Index Gymnastics}
\label{sec:sr-minkowski}

The geometry of flat spacetime is entirely encoded in a single object: the
metric.  Every interval, every dot product, every index-raising operation in
the codebase passes through it.  In curved spacetime the metric will vary from
point to point---that's the content of general relativity---but the algebraic
operations are identical.  Learning them here, where the metric is constant, is
like learning matrix multiplication with numbers before doing it with
functions.

\paragraph{Events and coordinates.}

An \emph{event} is a point in spacetime: a place and a time.  We label events
by four coordinates $x^\mu = (x^0, x^1, x^2, x^3) = (t, x, y, z)$, where the
index $\mu$ runs over $0, 1, 2, 3$.
\physnote{The superscript on $x^\mu$ is an \emph{index}, not an exponent.
Context always makes this clear.}
The coordinates themselves aren't physical---different observers may use
different labels for the same event---but the relationships between events
\emph{are} physical, and the metric is what makes them precise.

\paragraph{The Minkowski metric.}

The \emph{Minkowski metric} $\eta_{\mu\nu}$ assigns a notion of ``interval''
to pairs of nearby events.  In Cartesian coordinates, it takes the form
%
\begin{equation}
  \eta_{\mu\nu}
    = \operatorname{diag}(-1,\, +1,\, +1,\, +1) \,,
  \label{eq:minkowski-metric}
\end{equation}
%
so that the spacetime interval between two infinitesimally separated events is
%
\begin{equation}
  \d s^2
    = \eta_{\mu\nu}\,\d x^\mu\,\d x^\nu
    = -\d t^2 + \d x^2 + \d y^2 + \d z^2 \,.
  \label{eq:minkowski-interval}
\end{equation}
%
The crucial feature is the minus sign on $\d t^2$.  Without it, spacetime would
be four-dimensional Euclidean space and there would be no distinction between
time and space, no light cones, and no causality.  With it, we get the
three-way classification of directions that underpins all of relativistic
physics.
\warnnote{The other common convention is $(+{-}{-}{-})$.  We use $(-{+}{+}{+})$
throughout, following \cite{MTW1973}.  Formulae from sources using the opposite
convention will differ by an overall sign in certain places.}

We've used the \emph{Einstein summation convention}: any index appearing once
as a superscript and once as a subscript in a single term is summed over all
four values $0, 1, 2, 3$.  The expression $\eta_{\mu\nu}\,\d x^\mu\,\d
x^\nu$ therefore stands for the double sum
%
\begin{equation*}
  \eta_{\mu\nu}\,\d x^\mu\,\d x^\nu
    \;\equiv\;
    \sum_{\mu=0}^{3}\sum_{\nu=0}^{3}
    \eta_{\mu\nu}\,\d x^\mu\,\d x^\nu \,.
\end{equation*}
%
We will use this convention without further comment throughout the book.

\paragraph{Contravariant and covariant components.}

A \emph{four-vector} $V$ lives in the tangent space and has
\emph{contravariant} components $V^\mu$ (upper indices).  Its metric dual is a
covector with \emph{covariant} components $V_\mu$ (lower indices).  These are
geometrically distinct objects---the metric provides a canonical isomorphism
between them, so we pass freely between the two.  Lowering an index:
%
\begin{equation}
  V_\mu = \eta_{\mu\nu}\, V^\nu \,.
  \label{eq:index-lowering}
\end{equation}
%
Since $\eta_{\mu\nu}$ is diagonal with entries $(-1, +1, +1, +1)$, this
operation has a simple effect on the components: the spatial components are
unchanged, and the time component picks up a minus sign.  Explicitly,
$V_0 = -V^0$ while $V_i = V^i$ for $i = 1, 2, 3$.

Raising works with the \emph{inverse metric} $\eta^{\mu\nu}$, defined by the
condition $\eta^{\mu\alpha}\,\eta_{\alpha\nu} = \delta^\mu{}_\nu$.  For
Minkowski spacetime in Cartesian coordinates, the inverse happens to have the
same component values as the metric itself:
%
\begin{equation}
  \eta^{\mu\nu}
    = \operatorname{diag}(-1,\, +1,\, +1,\, +1) \,.
  \label{eq:minkowski-inverse}
\end{equation}
%
This is a coincidence of flat spacetime in Cartesian coordinates---for the
Schwarzschild or Kerr metrics, the inverse metric will have very different
components from the metric.  But the operation $V^\mu = \eta^{\mu\nu}\,V_\nu$
has the same structure in all cases: contract the metric (or its inverse) with
the vector.

Think of the metric as a translation dictionary between two different
languages.  Contravariant vectors and covariant covectors are different kinds of
geometric object: a velocity $\d x^\mu / \d\tau$ lives in the tangent space,
while a gradient $\partial_\mu \phi$ lives in the cotangent space.  They share
the same representation---an ordered list of four numbers---but they transform
differently and measure different things.  The metric provides an isomorphism
between the two spaces, converting one into the other freely.  This carries
through to curved spacetime, where the dictionary changes from point to point
but the two-space structure persists.
\xrefnote{In curved spacetime (\cref{ch:differential-geometry}),
$\eta_{\mu\nu}$ generalises to $\metric$ and the components become
functions of position.  The index operations look identical.}

\paragraph{The inner product and vector classification.}

The metric defines an inner product (indefinite, unlike the Euclidean case)
between any two four-vectors $A^\mu$ and $B^\mu$:
%
\begin{equation}
  A \cdot B
    = \eta_{\mu\nu}\, A^\mu\, B^\nu
    = -A^0 B^0 + A^1 B^1 + A^2 B^2 + A^3 B^3 \,.
  \label{eq:minkowski-inner-product}
\end{equation}
%
Unlike the Euclidean dot product, this can be negative, zero, or positive.  The
sign of a vector's ``squared norm'' $V^2 \equiv \eta_{\mu\nu}\,V^\mu\,V^\nu$
classifies it into one of three types:
%
\begin{itemize}
  \item \textbf{Timelike} ($V^2 < 0$): the vector points more into the future
    (or past) than into space.  Massive particles travel along timelike
    worldlines.
  \item \textbf{Null} ($V^2 = 0$): the vector lies exactly on the light cone.
    Photons travel along null geodesics---these are the trajectories that the
    ray tracer computes, millions of them per image.
  \item \textbf{Spacelike} ($V^2 > 0$): the vector points more into space
    than into time.  No massive particle or photon can travel along a spacelike
    direction.
\end{itemize}
%
This classification is \emph{invariant}: all observers agree on whether a
vector is timelike, null, or spacelike, regardless of their relative motion.
The metric makes this possible by providing an observer-independent notion of
``norm''.
\physnote{The null condition $V^2 = 0$ doesn't mean the vector is zero---it
means the time and space contributions cancel exactly.  A photon moves through
space precisely as fast as it moves through time.}

\begin{figure}[htbp]
  \centering
  % minkowski-lightcone.tex — Light cone and causal vector classes in 1+1 Minkowski space
% Inputable TikZ fragment for fig:minkowski-lightcone
% Requires: tikz with calc, arrows.meta (loaded by ehbook.cls)
%
\begin{tikzpicture}[
    >={Stealth[length=5pt]},
    font=\footnotesize,
  ]

  \def\axlen{2.8}     % axis half-length
  \def\cone{2.4}      % light-cone extent along each ray

  % ── Shaded future / past regions ──────────────────────────────────
  \fill[EHAccretionGold, fill opacity=0.08]
    (0,0) -- (-\cone,\cone) -- (\cone,\cone) -- cycle;
  \fill[EHJetBlue, fill opacity=0.06]
    (0,0) -- (-\cone,-\cone) -- (\cone,-\cone) -- cycle;

  % Region labels
  \node[font=\scriptsize, text=EHMarginGray] at (0, 1.7) {future};
  \node[font=\scriptsize, text=EHMarginGray] at (0,-1.7) {past};

  % ── Light-cone edges ──────────────────────────────────────────────
  \draw[EHMarginGray, thin] (-\cone,\cone)  -- (0,0) -- (\cone,\cone);
  \draw[EHMarginGray, thin] (-\cone,-\cone) -- (0,0) -- (\cone,-\cone);

  % ── Axes ──────────────────────────────────────────────────────────
  \draw[->, thick, EHMarginGray] (-\axlen,0) -- (\axlen,0)
    node[right, font=\small] {$x$};
  \draw[->, thick, EHMarginGray] (0,-\axlen) -- (0,\axlen)
    node[above, font=\small] {$t$};

  % ── Timelike vector (gold, inside future cone) ────────────────────
  \draw[->, very thick, EHAccretionGold]
    (0,0) -- (0.45,1.4)
    node[left=2pt, text=EHAccretionGold!85!black] {timelike};

  % ── Spacelike vector (blue, outside cone) ─────────────────────────
  \draw[->, very thick, EHJetBlue]
    (0,0) -- (1.9,0.4)
    node[right=2pt, text=EHJetBlue!85!black] {spacelike};

  % ── Null vector (crimson, along cone edge) ────────────────────────
  \draw[->, very thick, EHHorizonCrimson]
    (0,0) -- (1.55,1.55)
    node[right=2pt, text=EHHorizonCrimson!85!black] {null};

\end{tikzpicture}

  \caption[Causal classes of four-vectors in Minkowski space]{Spacetime diagram of 1+1 Minkowski space, showing the three causal
    classes of four-vectors.  Timelike vectors (gold) lie inside the light cone
    and have negative norm $\eta_{\mu\nu} V^\mu V^\nu < 0$; spacelike vectors
    (blue) lie outside and have positive norm; null vectors (crimson) lie
    exactly on the cone with vanishing norm.  The photon trajectories computed
    by the ray tracer are always null---they live on this cone and its
    curved-spacetime generalisation.}
  \label{fig:minkowski-lightcone}
\end{figure}

\paragraph{A worked example.}

Consider a photon travelling in the $x$-direction.  Over a coordinate time
interval $\Delta t$, its displacement four-vector is
%
\begin{equation}
  \Delta x^\mu = (\Delta t,\, \Delta t,\, 0,\, 0) \,,
  \label{eq:photon-displacement}
\end{equation}
%
where we use $c = 1$.  Contract with the metric:
%
\begin{equation}
  \Delta s^2
    = \eta_{\mu\nu}\,\Delta x^\mu\, \Delta x^\nu
    = -(\Delta t)^2 + (\Delta t)^2 + 0 + 0
    = 0 \,.
  \label{eq:photon-interval}
\end{equation}
%
The interval vanishes: the displacement is null, exactly as required for a
photon.  This isn't a derived result---it's the \emph{definition} of what it
means to travel at the speed of light.  The entire ray tracer is built on
propagating this constraint through curved spacetime, where the metric changes
from $\eta_{\mu\nu}$ to $\metric$ and the algebra becomes a system of
differential equations.

For a spacelike example, take two events at the same coordinate time separated
by a distance $d$: $\Delta x^\mu = (0,\, d,\, 0,\, 0)$.  Then $\Delta s^2 =
d^2 > 0$---spacelike, as expected for a purely spatial separation.

Now consider a massive particle moving in the $x$-direction at speed $v < 1$.
Over a coordinate time $\Delta t$, the displacement is $\Delta x^\mu = (\Delta
t,\, v\,\Delta t,\, 0,\, 0)$, giving
%
\begin{equation}
  \Delta s^2
    = -(1 - v^2)\,(\Delta t)^2
    < 0 \,.
  \label{eq:timelike-interval}
\end{equation}
%
The interval is negative: the displacement is timelike.  For any timelike
displacement ($\Delta s^2 < 0$), the quantity
$\Delta\tau = \sqrt{-\Delta s^2}$ is the proper time elapsed on the particle's
own clock---less than the coordinate time $\Delta t$ for any $v > 0$, and
equal to it in the rest frame.

\paragraph{Looking ahead.}

The single minus sign in $\eta_{\mu\nu}$ is the source of all causal
structure: it splits spacetime into timelike, null, and spacelike directions,
gives us light cones, and makes the invariant interval possible.  Without it,
there is no distinction between time and space---and no ray tracing.

Everything we've done in this section generalises directly.  When we move to
curved spacetime in \cref{ch:differential-geometry}, the constant matrix
$\eta_{\mu\nu}$ is replaced by a position-dependent tensor field $\metric$,
the interval becomes $\d s^2 = \metric\,\d x^\mu\,\d x^\nu$, and the index
operations acquire exactly the same form but with $\metric$ and $\inversemetric$
in place of $\eta_{\mu\nu}$ and $\eta^{\mu\nu}$.  The classification of
vectors into timelike, null, and spacelike is unchanged---only the specific
components of the metric are different.  This is why we invest in learning the
index machinery here: it's the same machinery everywhere, and once it's
automatic, the physics of curved spacetime is free to take centre stage.


% sec:sr-lorentz — Lorentz Transformations, Boosts, and Observer Tetrads
\section{Lorentz Transformations, Boosts, and Observer Tetrads}
\label{sec:sr-lorentz}

Two physicists in relative motion measure the same pair of events and compute
the spacetime interval.  Both must get the same answer---that is the content of
the invariance principle from \cref{sec:sr-minkowski}.  But they use different
coordinates, so their individual numbers for $\Delta t$ and $\Delta x$ differ.
Which coordinate changes are compatible with this requirement?  The answer is
the Lorentz group, and its most physically important elements---boosts---are
the transformations that connect observers in relative motion.  The tetrad
construction at the end of this section packages an observer's local frame into
four vectors that generalise directly to curved spacetime, where they define the
ray tracer's camera.

\paragraph{The Lorentz group.}

A \emph{Lorentz transformation} is a linear map $x'^\mu = \Lambda^\mu{}_\nu\,
x^\nu$ that preserves the Minkowski interval.  The preservation condition is
that for any pair of events,
%
\begin{equation}
  \eta_{\mu\nu}\,\d x'^\mu\,\d x'^\nu
    = \eta_{\mu\nu}\,\d x^\mu\,\d x^\nu \,.
  \label{eq:interval-invariance}
\end{equation}
%
Substituting $\d x'^\mu = \Lambda^\mu{}_\alpha\,\d x^\alpha$ and requiring this
to hold for arbitrary $\d x^\alpha$, we obtain the defining condition
%
\begin{equation}
  \eta_{\mu\nu}\,\Lambda^\mu{}_\alpha\,\Lambda^\nu{}_\beta
    = \eta_{\alpha\beta} \,.
  \label{eq:lorentz-condition}
\end{equation}
%
Read this as a constraint on the matrix $\Lambda$: the Minkowski metric is
``sandwiched'' between two copies of $\Lambda$ and must reproduce itself.  This
is the hyperbolic analogue of the orthogonality condition $R^T R = I$ in
Euclidean geometry---there, the identity matrix plays the role that $\eta$ plays
here.  In matrix notation, \cref{eq:lorentz-condition} reads
%
\begin{equation}
  \Lambda^{T}\!\eta\,\Lambda = \eta \,.
  \label{eq:lorentz-matrix}
\end{equation}
%
Any $4 \times 4$ real matrix satisfying this condition belongs to the
\emph{Lorentz group} $O(1,3)$.
\histnote{The group is named after Hendrik Lorentz, who derived the
transformations in 1904.  Einstein's insight was that they express a physical
symmetry, not merely a mathematical convenience.}

Taking the determinant of both sides of \cref{eq:lorentz-matrix} gives
$(\det\Lambda)^2 = 1$, so $\det\Lambda = \pm 1$.  Setting $\alpha = \beta = 0$
in \cref{eq:lorentz-condition} gives
$-(\Lambda^0{}_0)^2 + \sum_i (\Lambda^i{}_0)^2 = -1$; since the spatial sum
$\sum_i (\Lambda^i{}_0)^2$ is non-negative, it follows that
$(\Lambda^0{}_0)^2 \geq 1$.  The \emph{proper orthochronous} Lorentz group
$SO^+(1,3)$ consists of those elements with $\det\Lambda = +1$ and
$\Lambda^0{}_0 \geq 1$---these are the physically realisable transformations
that preserve spatial orientation and the direction of time.  We restrict
attention to this subgroup.

\paragraph{Boosts.}

The simplest proper orthochronous Lorentz transformations that mix space and
time are \emph{boosts}.  Consider two observers: $S$, at rest, and $S'$, moving
with velocity $v$ in the $x$-direction relative to $S$ (with $c = 1$, so
$v \in [0,1)$).  The transformation connecting their coordinates must satisfy
\cref{eq:lorentz-condition}, leave $y$ and $z$ unchanged, and reduce to the
identity at $v = 0$.  The unique solution is
%
\begin{equation}
  \Lambda^\mu{}_\nu
    = \begin{pmatrix}
        \gamma   & -v\gamma & 0 & 0 \\
        -v\gamma & \gamma   & 0 & 0 \\
        0        & 0        & 1 & 0 \\
        0        & 0        & 0 & 1
      \end{pmatrix} ,
  \label{eq:boost-matrix}
\end{equation}
%
where the \emph{Lorentz factor} is
%
\begin{equation}
  \gamma = \frac{1}{\sqrt{1 - v^2}} \,.
  \label{eq:lorentz-factor}
\end{equation}
%
This is the passive transformation giving $S'$ coordinates in terms of $S$
coordinates, with $S'$ moving at velocity $+v$ along $x$.  At $v = 0$ the
matrix is the identity; as $v \to 1$ the off-diagonal entries diverge,
reflecting the impossibility of boosting to the speed of light.
\warnnote{The signs in \cref{eq:boost-matrix} depend on the signature
convention.  In $(+{-}{-}{-})$ sources, the off-diagonal entries have
the opposite sign.}

It is worth verifying \cref{eq:lorentz-condition} explicitly.  The $(0,0)$
component gives $-\gamma^2 + v^2\gamma^2 = -\gamma^2(1 - v^2) = -1$, which
holds by definition of $\gamma$.  The $(0,1)$ component gives $v\gamma^2 -
v\gamma^2 = 0$, confirming the off-diagonal entries of $\eta$ vanish.  The
spatial block reduces to the identity since $\Lambda$ acts trivially on $y$ and
$z$.

\paragraph{Rapidity.}

An alternative parametrisation replaces $v$ with the \emph{rapidity} $\phi$,
defined by $v = \tanh\phi$ (equivalently, $\gamma = \cosh\phi$ and $v\gamma =
\sinh\phi$).  The boost matrix becomes
%
\begin{equation}
  \Lambda^\mu{}_\nu
    = \begin{pmatrix}
        \cosh\phi  & -\sinh\phi & 0 & 0 \\
        -\sinh\phi & \cosh\phi  & 0 & 0 \\
        0          & 0          & 1 & 0 \\
        0          & 0          & 0 & 1
      \end{pmatrix} .
  \label{eq:boost-rapidity}
\end{equation}
%
The payoff is that rapidities add: two successive collinear boosts with
rapidities $\phi_1$ and $\phi_2$ compose to a single boost with rapidity
$\phi_1 + \phi_2$.  This follows from the hyperbolic addition formula:
$\cosh(\phi_1 + \phi_2) = \cosh\phi_1\cosh\phi_2 + \sinh\phi_1\sinh\phi_2$,
which is exactly the $(0,0)$ entry of the matrix product
$\Lambda(\phi_1)\,\Lambda(\phi_2)$.  Contrast this with the relativistic
velocity-addition formula $w = (v_1 + v_2)/(1 + v_1 v_2)$---velocities compose
nonlinearly, but their corresponding rapidities simply sum.  For non-collinear
boosts, the composition also produces a spatial rotation (the Wigner rotation).
\physnote{Rapidity is the hyperbolic analogue of the rotation angle.  Rotations
compose by adding angles; boosts compose by adding rapidities.}

\paragraph{Observer tetrads.}

A Lorentz transformation connects two coordinate systems.  A \emph{tetrad}
goes further: it is an observer's personal ruler-and-clock kit planted at a
single point in spacetime.  One vector says ``this is my future direction,'' and
three mutually perpendicular vectors say ``these are my spatial axes.''  The
formal requirement is that these four vectors reproduce the Minkowski metric
when measured against the background.

An \emph{orthonormal tetrad} is a set of four vectors
$e^\mu{}_{(a)}$, with $a = 0, 1, 2, 3$, satisfying
%
\begin{equation}
  \eta_{\mu\nu}\, e^\mu{}_{(a)}\, e^\nu{}_{(b)}
    = \eta_{ab} \,.
  \label{eq:tetrad-orthonormality}
\end{equation}
%
The parenthesised index $(a)$ is a \emph{frame index}: it labels which tetrad
vector we are referring to, not a spacetime direction.  The notation
$e^\mu{}_{(a)}$ carries one contravariant spacetime index $\mu$ (giving the
vector's components in coordinates) and one covariant frame index $(a)$
(labelling which basis vector); the inverse tetrad $e^{(a)}{}_\mu$ reverses
these roles.  The orthonormality condition says that the timelike vector
$e^\mu{}_{(0)}$ has norm $-1$ and is orthogonal to the three spacelike vectors
$e^\mu{}_{(i)}$, each of norm $+1$.
\warnnote{Tetrad indices $(a), (b), \ldots$ live in a \emph{local} Minkowski
frame.  They are not raised or lowered with the spacetime metric---they use
$\eta_{ab}$ directly.}

For an observer at rest in the standard coordinates, the tetrad is trivial:
%
\begin{equation}
  \begin{aligned}
    e^\mu{}_{(0)} &= (1,0,0,0), \quad &
    e^\mu{}_{(1)} &= (0,1,0,0), \\
    e^\mu{}_{(2)} &= (0,0,1,0), \quad &
    e^\mu{}_{(3)} &= (0,0,0,1) \,.
  \end{aligned}
  \label{eq:tetrad-rest}
\end{equation}
%
This is the identity matrix---unsurprising, since the observer's frame
coincides with the coordinate frame.  For an observer moving with velocity
$v$ in the $x$-direction, we apply the inverse boost (replacing $v \to -v$ in
\cref{eq:boost-matrix}) to the rest-frame tetrad vectors, transforming them
from the moving observer's frame into $S$ coordinates:
%
\begin{equation}
  e^\mu{}_{(0)} = (\gamma,\, v\gamma,\, 0,\, 0), \quad
  e^\mu{}_{(1)} = (v\gamma,\, \gamma,\, 0,\, 0) \,.
  \label{eq:tetrad-boosted}
\end{equation}
%
The vectors $e^\mu{}_{(2)}$ and $e^\mu{}_{(3)}$ are unchanged.  The timelike
tetrad vector $e^\mu{}_{(0)} = \fourvel$ is the observer's
\emph{four-velocity}---this is not a coincidence but the formal definition of
``the direction the observer moves through spacetime.''  The tetrad is a
coordinate-free encoding of the observer's state of motion: it captures both
their velocity (via $e^\mu{}_{(0)}$) and their spatial orientation (via
$e^\mu{}_{(i)}$), which spans the observer's simultaneity surface and defines
the directions they call ``right,'' ``up,'' and ``forward.''

\paragraph{A worked example: camera frame at $v = 0.6$.}

Consider an observer (a camera, in our context) moving at $v = 0.6$ in the
$x$-direction.  The Lorentz factor is $\gamma = 1/\sqrt{1 - 0.36} = 1.25$.
The timelike tetrad vector is
$e^\mu{}_{(0)} = (1.25,\, 0.75,\, 0,\, 0)$,
and we can verify its norm:
$\eta_{\mu\nu}\,e^\mu{}_{(0)}\,e^\nu{}_{(0)} = -(1.25)^2 + (0.75)^2 =
-1.5625 + 0.5625 = -1$.  The spatial vector
$e^\mu{}_{(1)} = (0.75,\, 1.25,\, 0,\, 0)$ has norm
$-(0.75)^2 + (1.25)^2 = -0.5625 + 1.5625 = +1$, and the inner product
$\eta_{\mu\nu}\,e^\mu{}_{(0)}\,e^\nu{}_{(1)} = -(1.25)(0.75) + (0.75)(1.25)
= 0$.  The tetrad is indeed orthonormal.

\begin{figure}[htbp]
  \centering
  % boosted-tetrad.tex — Boosted tetrad vectors in 1+1 Minkowski space
% Inputable TikZ fragment for fig:boosted-tetrad
% Requires: tikz with calc, arrows.meta (loaded by ehbook.cls)
%
% Boost v = 0.6 → γ = 1.25
% e_{(0)} = (γ, γv) = (1.25, 0.75)   — timelike, tilted from t-axis
% e_{(1)} = (γv, γ) = (0.75, 1.25)   — spacelike, tilted from x-axis
% Tilt angle: arctan(0.6) ≈ 31°
%
\begin{tikzpicture}[
    >={Stealth[length=5pt]},
    font=\footnotesize,
  ]

  \def\axlen{2.8}     % axis half-length
  \def\cone{2.4}      % light-cone extent along each ray
  \def\vecscale{1.4}  % scale factor for tetrad vectors

  % ── Light-cone shading ────────────────────────────────────────────
  \fill[EHMarginGray, fill opacity=0.05]
    (0,0) -- (-\cone,\cone) -- (\cone,\cone) -- cycle;
  \fill[EHMarginGray, fill opacity=0.05]
    (0,0) -- (-\cone,-\cone) -- (\cone,-\cone) -- cycle;

  % ── Light-cone edges ──────────────────────────────────────────────
  \draw[EHMarginGray, thin] (-\cone,\cone)  -- (0,0) -- (\cone,\cone);
  \draw[EHMarginGray, thin] (-\cone,-\cone) -- (0,0) -- (\cone,-\cone);

  % ── Lab frame axes (S) ───────────────────────────────────────────
  \draw[->, thick, EHMarginGray] (-\axlen,0) -- (\axlen,0)
    node[right, font=\small] {$x$};
  \draw[->, thick, EHMarginGray] (0,-\axlen) -- (0,\axlen)
    node[above, font=\small] {$t$};

  % ── Boosted axes (dashed, to show tilted coordinate grid) ────────
  % t'-axis along e_{(0)} direction: slope = γ/(γv) = 1/v = 1/0.6
  \draw[EHMarginGray, thin, densely dashed]
    (0,0) -- ({0.6*\cone},\cone)
    node[right, font=\scriptsize, text=EHMarginGray] {$t'$};
  % x'-axis along e_{(1)} direction: slope = γv/γ = v = 0.6
  \draw[EHMarginGray, thin, densely dashed]
    (0,0) -- (\cone,{0.6*\cone})
    node[above, font=\scriptsize, text=EHMarginGray] {$x'$};

  % ── Timelike tetrad vector e_{(0)} (gold) ────────────────────────
  % e_{(0)} = (γ, γv) = (1.25, 0.75), plotted as (x,t) = (0.75, 1.25)
  \draw[->, very thick, EHAccretionGold]
    (0,0) -- ({0.75*\vecscale},{1.25*\vecscale})
    node[left=2pt, text=EHAccretionGold!85!black]
      {$e^\mu{}_{(0)}$};

  % ── Spacelike tetrad vector e_{(1)} (blue) ───────────────────────
  % e_{(1)} = (γv, γ) = (0.75, 1.25), plotted as (x,t) = (1.25, 0.75)
  \draw[->, very thick, EHJetBlue]
    (0,0) -- ({1.25*\vecscale},{0.75*\vecscale})
    node[right=2pt, text=EHJetBlue!85!black]
      {$e^\mu{}_{(1)}$};

\end{tikzpicture}

  \caption[Boosted tetrad vectors]{Tetrad vectors for an observer moving at $v = 0.6$ in the
    $x$-direction.  The timelike vector $e^\mu{}_{(0)}$ (gold) is the
    observer's four-velocity, tilted toward the $x$-axis by the boost.  The
    spatial vector $e^\mu{}_{(1)}$ (blue) is the observer's $x'$-direction,
    tilted symmetrically toward the $t$-axis.  Together they satisfy the
    orthonormality condition
    $\eta_{\mu\nu}\,e^\mu{}_{(a)}\,e^\nu{}_{(b)} = \eta_{ab}$.  In curved
    spacetime, the same construction defines the camera's local reference
    frame at any point.}
  \label{fig:boosted-tetrad}
\end{figure}

\paragraph{From flat to curved.}
\coderef{Spacetime}{Minkowski}

In the codebase, the Minkowski metric is an instance of the \texttt{Spacetime}
typeclass, with \texttt{metric4} returning \texttt{sym4Diag (-1) 1 1 1} and all
Christoffel symbols set to zero.  The same typeclass provides Schwarzschild and
Kerr metrics, and the camera's tetrad construction works identically for all of
them---only the orthonormality condition changes from $\eta_{\mu\nu}$ to
$\metric$.
\implnote{The \texttt{Spacetime} typeclass abstracts over the specific metric,
so the tetrad and ray-tracing code is metric-agnostic.}

In curved spacetime, the tetrad condition generalises to
%
\begin{equation}
  g_{\mu\nu}\,e^\mu{}_{(a)}\,e^\nu{}_{(b)} = \eta_{ab} \,.
  \label{eq:tetrad-curved}
\end{equation}
%
The metric on the left is the curved-spacetime metric, but the right-hand side
remains the flat Minkowski matrix.  The tetrad defines a locally Minkowski frame
at each point, regardless of the background curvature.  The camera model in
\cref{ch:camera-rendering} uses exactly this construction to define the
observer's ``sky''---the directions from which photons arrive---at any point in
any spacetime.
\xrefnote{The tetrad reappears in the camera model of
\cref{ch:camera-rendering}, where it defines the mapping from pixel
coordinates to initial photon momenta.}


% sec:sr-causal — Light Cones, Causal Structure, and Null Vectors
\section{Light Cones, Causal Structure, and Null Vectors}
\label{sec:sr-causal}

A ray tracer needs to answer one question for every pixel: given a camera at
some event in spacetime, which direction does this pixel correspond to, and
where does a photon from that direction originate?  The answer lives in the set
of null vectors at the camera's location---the directions along which light can
travel.  This set has a rich geometric structure: it forms a cone in the tangent
space, and the spatial directions on that cone are in one-to-one correspondence
with points on the observer's sky.  This section makes that correspondence
precise.

\paragraph{The light cone.}

At any event $p$ in Minkowski spacetime, the \emph{light cone} is the set of
all null directions---nonzero vectors $k^\mu$ satisfying
%
\begin{equation}
  \eta_{\mu\nu}\, k^\mu\, k^\nu = 0 \,.
  \label{eq:null-condition}
\end{equation}
%
In \cref{sec:sr-minkowski} we classified individual vectors as timelike, null,
or spacelike.  The light cone is the \emph{boundary} that separates the
timelike region from the spacelike region: away from the vertex $k^\mu = 0$,
it is a three-dimensional hypersurface in the four-dimensional tangent space
at $p$.

Expanding \cref{eq:null-condition} in Cartesian coordinates gives
%
\begin{equation}
  -(k^0)^2 + (k^1)^2 + (k^2)^2 + (k^3)^2 = 0 \,,
  \label{eq:null-expanded}
\end{equation}
%
which rearranges to $(k^0)^2 = \abs{\mathbf{k}}^2$, where $\abs{\mathbf{k}}^2
= (k^1)^2 + (k^2)^2 + (k^3)^2$ is the squared Euclidean norm of the spatial
part.  This has two solutions: $k^0 = +\abs{\mathbf{k}}$ (future-directed) and
$k^0 = -\abs{\mathbf{k}}$ (past-directed).  The light cone therefore splits
into two halves:
%
\begin{itemize}
  \item The \textbf{future light cone}: all null vectors with $k^0 > 0$.
    These point toward the future---photons emitted from $p$ travel along
    these directions.
  \item The \textbf{past light cone}: all null vectors with $k^0 < 0$.
    These point toward the past---photons arriving at $p$ came from
    these directions.
\end{itemize}
%
The ray tracer works with the past light cone: it traces photons backward in
time from the camera to their sources, because that is far more efficient than
tracing photons forward from every source and hoping some reach the camera.
\physnote{Backward ray tracing is standard in computer graphics
\cite{Marschner2016}; the relativistic version merely replaces straight lines
with null geodesics.}

\begin{figure}[htbp]
  \centering
  % light-cone-causal.tex — 2+1 light cone with causal regions
% Inputable TikZ fragment for fig:light-cone-causal
% Requires: tikz with calc, arrows.meta, decorations.markings (loaded by ehbook.cls)
%
\begin{tikzpicture}[
    >={Stealth[length=5pt]},
    font=\footnotesize,
  ]

  \def\h{2.2}     % cone height
  \def\rx{1.65}   % ellipse semi-major (x direction)
  \def\ry{0.55}   % ellipse semi-minor (perspective depth)

  % Tangent points: lines from apex (0,0) tangent to ellipse at height h.
  %   tx = rx·sqrt(h²-ry²)/h,  ty = (h²-ry²)/h,  tangle = acos(tx/rx)
  \pgfmathsetmacro{\tx}{\rx*sqrt(\h*\h - \ry*\ry)/\h}
  \pgfmathsetmacro{\ty}{(\h*\h - \ry*\ry)/\h}
  \pgfmathsetmacro{\tangle}{acos(\tx/\rx)}

  % ── Past cone ─────────────────────────────────────────────────────

  % Back arc of past ellipse (dashed)
  \draw[EHMarginGray, thin, densely dashed]
    (-\tx,-\ty) arc ({180-\tangle}:{360+\tangle}:{\rx} and {\ry});

  % Past cone fill (very light)
  \fill[EHJetBlue, fill opacity=0.05]
    (0,0) -- (\tx,-\ty) arc ({\tangle}:{180-\tangle}:{\rx} and {\ry}) -- cycle;

  % Past cone surface lines
  \draw[EHMarginGray, thin] (0,0) -- (-\tx,-\ty);
  \draw[EHMarginGray, thin] (0,0) -- ( \tx,-\ty);

  % Front arc of past ellipse
  \draw[EHMarginGray, thin]
    (\tx,-\ty) arc ({\tangle}:{180-\tangle}:{\rx} and {\ry});

  % ── Axes (behind future cone) ─────────────────────────────────────

  \draw[->, thick, EHMarginGray] (-2.6,0) -- (2.6,0)
    node[right, font=\small] {$x$};
  \draw[->, thick, EHMarginGray] (0,0) -- (-0.85,-0.5)
    node[below left, font=\small] {$y$};

  % ── Future cone ───────────────────────────────────────────────────

  % Back arc of future ellipse (dashed)
  \draw[EHHorizonCrimson!50, thin, densely dashed]
    (\tx,\ty) arc ({-\tangle}:{180+\tangle}:{\rx} and {\ry});

  % Future cone interior fill — I^+(p)
  \fill[EHAccretionGold, fill opacity=0.10]
    (0,0) -- (-\tx,\ty)
      arc ({180+\tangle}:{360-\tangle}:{\rx} and {\ry}) -- cycle;

  % Cone surface lines — boundary of J^+(p)
  \draw[EHHorizonCrimson, thick] (0,0) -- (-\tx,\ty);
  \draw[EHHorizonCrimson, thick] (0,0) -- ( \tx,\ty);

  % Front arc of future ellipse
  \draw[EHHorizonCrimson, thick]
    (-\tx,\ty) arc ({180+\tangle}:{360-\tangle}:{\rx} and {\ry});

  % ── t-axis (drawn last so it's on top) ────────────────────────────

  \draw[->, thick, EHMarginGray] (0,-2.8) -- (0,2.8)
    node[above, font=\small] {$t$};

  % ── Event p ────────────────────────────────────────────────────────

  \fill (0,0) circle (1.5pt);
  \node[below right=1pt, font=\small] at (0,0) {$p$};

  % ── Timelike worldline (gold, through cone interior) ──────────────

  \draw[very thick, EHAccretionGold,
        decoration={markings, mark=at position 0.72 with {\arrow{>}}},
        postaction={decorate}]
    (0.15,-1.5) .. controls (0.05,-0.5) and (-0.1,0.5) .. (-0.25,1.8);

  % ── Spacelike separation (blue, outside cone) ─────────────────────

  \draw[->, very thick, EHJetBlue]
    (0,0) -- (2.1,0.2);

  % ── Labels ─────────────────────────────────────────────────────────

  \node[text=EHAccretionGold!85!black] at (0.4,1.15) {$I^+(p)$};
  \node[text=EHHorizonCrimson!85!black, right=1pt] at (\tx,\ty) {$J^+(p)$};
  \node[text=EHJetBlue!85!black, above=1pt] at (2.1,0.2) {elsewhere};

\end{tikzpicture}

  \caption[The light cone at an event $p$ in Minkowski spacetime]{The light cone at an event $p$ in Minkowski spacetime.  The
    chronological future $I^+(p)$ (gold interior) contains all events
    reachable by timelike curves---massive particles.  The causal future
    $J^+(p)$ adds the cone boundary itself (crimson), reached by null
    geodesics---photons.  Events outside both cones are spacelike-separated
    from $p$: no signal can connect them.  The ray tracer works with the past
    light cone, tracing photons backward from the camera to their sources.}
  \label{fig:light-cone-causal}
\end{figure}

\paragraph{Causal structure.}

The light cone at an event $p$ divides the tangent space---and, by extension,
the rest of spacetime---into causally distinct regions.  The \emph{chronological
future} $I^+(p)$ is the set of events reachable from $p$ by timelike curves
(massive particles).  The \emph{causal future} $J^+(p)$ includes both timelike
and null curves---it adds the events connected to $p$ by light signals.  The
boundary of $J^+(p)$ is precisely the future light cone of $p$.  Analogous
definitions hold for the past: $I^-(p)$ and $J^-(p)$.

Events outside $J^+(p) \cup J^-(p)$ are \emph{causally disconnected} from $p$:
no signal, massive or massless, can travel between them.  The separation
between $p$ and any such event is spacelike.  This is the geometric content of
the statement that nothing travels faster than light---it is not a speed limit
imposed from outside, but a consequence of the cone structure built into the
metric.
\histnote{Penrose's conformal diagrams \cite{Penrose1963} make the global
causal structure of a spacetime visible at a glance.  We will not need them
until \cref{ch:schwarzschild}, where the causal structure of the black hole
interior becomes central.}

\paragraph{Null vectors and the observer's sky.}

The null condition \cref{eq:null-condition} reduces a vector from four
components down to three free parameters.  But there is a further redundancy: if $k^\mu$ is
null, so is $\lambda k^\mu$ for any $\lambda > 0$.  Rescaling a photon's
four-momentum changes its energy but not its direction.  The physically
distinct null directions therefore form a \emph{two}-parameter family---they
are labelled by a point on a two-sphere.

To see this explicitly, write a future-directed null vector as
%
\begin{equation}
  k^\mu = E\,(1,\, \sin\theta\cos\varphi,\,
               \sin\theta\sin\varphi,\, \cos\theta) \,,
  \label{eq:null-parametrisation}
\end{equation}
%
where $E > 0$ is the photon energy (the freedom we quotient out) and $(\theta,
\varphi)$ are the standard spherical angles on $S^2$.  The two angles label the
directions on the observer's \emph{celestial sphere}---the sky.  Each pixel in
a rendered image corresponds to a specific $(\theta, \varphi)$, and therefore
to a specific null direction $k^\mu$.  Substituting into
\cref{eq:null-condition} confirms the null condition: $-E^2 + E^2(\sin^2\theta
\cos^2\varphi + \sin^2\theta\sin^2\varphi + \cos^2\theta) = -E^2 + E^2 = 0$.

This is the starting point for every ray trace: pick a pixel, compute
the null vector, and propagate it through spacetime.
\xrefnote{The camera model in \cref{ch:camera-rendering} maps pixel
coordinates $(i, j)$ to angles $(\theta, \varphi)$ using the observer's
tetrad from \cref{sec:sr-lorentz}.}

\paragraph{The null covector and momentum.}

Lowering the index on a null vector yields a covector, the photon's
four-momentum $p_\mu$:
%
\begin{equation}
  p_\mu = \eta_{\mu\nu}\, k^\nu \,.
  \label{eq:photon-momentum}
\end{equation}
%
With signature $(-,+,+,+)$, the time component changes sign:
$p_0 = -k^0 = -E$, while the spatial components are unchanged in Cartesian
coordinates: $p_i = k^i$.  The null condition expressed in terms of the
covector is $\eta^{\mu\nu}\,p_\mu\,p_\nu = 0$, which reads
%
\begin{equation}
  -(p_0)^2 + (p_1)^2 + (p_2)^2 + (p_3)^2 = 0 \,,
  \label{eq:null-covector}
\end{equation}
%
or equivalently $p_0 = -\abs{\mathbf{p}}$, where the sign reflects our choice
that $p_0 < 0$ for a future-directed photon.  The covector form is the natural
variable for the ray tracer: the time component $p_0 = -E$ is minus the photon
energy as measured by a static observer, and the spatial components encode its
direction of travel.
\physnote{The minus sign on $p_0$ is a perennial source of sign errors.  In our
$(-,+,+,+)$ convention, a physical photon has $p_0 < 0$ and $E = -p_0 > 0$.}

The Hamiltonian ray tracer integrates the geodesic equations in terms of
$p_\mu$ because the covariant components have simpler evolution equations in
the presence of symmetry---along a geodesic, a Killing vector $\xi^\mu$ gives a
conserved quantity $p_\mu\,\xi^\mu$ directly, without further metric
contraction.

\paragraph{A worked example: constructing a photon's initial data.}

Suppose a camera at rest in Minkowski spacetime (so its tetrad is the identity)
needs to trace a ray toward the direction $(\theta, \varphi) = (60°, 45°)$.
The relevant trigonometric values are $\sin 60° = \sqrt{3}/2$,
$\cos 60° = 1/2$, $\cos 45° = \sin 45° = 1/\sqrt{2}$, and
$\sin 60°\cos 45° = \sqrt{3}/(2\sqrt{2}) = \sqrt{6}/4$.  The initial null
vector is therefore
%
\begin{equation}
  k^\mu = E\bigl(1,\, \tfrac{\sqrt{6}}{4},\,
                     \tfrac{\sqrt{6}}{4},\, \tfrac{1}{2}\bigr) \,.
  \label{eq:photon-example}
\end{equation}
%
The covector is $p_\mu = (-E,\, E\sqrt{6}/4,\,
E\sqrt{6}/4,\, E/2)$.  We can verify the null condition:
$\eta^{\mu\nu} p_\mu p_\nu = -E^2 + 6E^2/16 + 6E^2/16 + E^2/4
= -E^2 + 3E^2/8 + 3E^2/8 + 2E^2/8 = -E^2 + E^2 = 0$.

The energy $E$ drops out of the geodesic equations for null rays (the
trajectory depends only on direction, not on energy), so we can set $E = 1$
without loss of generality.  This gives the concrete initial four-momentum
$p_\mu = (-1,\, \sqrt{6}/4,\, \sqrt{6}/4,\, 1/2)$ that the integrator takes as
input.
\warnnote{Setting $E = 1$ is a choice of affine parametrisation, not a physical
statement about the photon's energy.  Observables like redshift involve
\emph{ratios} of energies measured by different observers, so the absolute
normalisation is irrelevant.}

\paragraph{From flat to curved: the sky in any spacetime.}

In Minkowski spacetime, the null parametrisation \cref{eq:null-parametrisation}
uses the coordinate axes as the reference frame.  For a boosted observer, we
replace the coordinate basis with the observer's tetrad from
\cref{sec:sr-lorentz}: the spatial directions
$(\sin\theta\cos\varphi,\, \sin\theta\sin\varphi,\, \cos\theta)$
are expanded in the tetrad vectors $e^\mu{}_{(1)}, e^\mu{}_{(2)},
e^\mu{}_{(3)}$ instead of the coordinate vectors, and the time direction uses
$e^\mu{}_{(0)}$.  The null vector for a pixel at $(\theta, \varphi)$ as
measured by observer $\mathcal{O}$ becomes
%
\begin{equation}
  k^\mu = E\Bigl(e^\mu{}_{(0)}
    + \sin\theta\cos\varphi\; e^\mu{}_{(1)}
    + \sin\theta\sin\varphi\; e^\mu{}_{(2)}
    + \cos\theta\; e^\mu{}_{(3)}\Bigr) \,.
  \label{eq:null-tetrad}
\end{equation}
%
This expression is valid in \emph{any} spacetime: flat or curved, static or
dynamic.  The tetrad encodes the observer's state of motion; the angles encode
the pixel direction; the metric enters only through the tetrad orthonormality
condition.  To confirm that \cref{eq:null-tetrad} is indeed null, substitute
into $\eta_{\mu\nu}\,k^\mu\,k^\nu$ and use the orthonormality condition
$\eta_{\mu\nu}\,e^\mu{}_{(a)}\,e^\nu{}_{(b)} = \eta_{ab}$: the cross terms
vanish, the timelike leg contributes $-E^2$, and the three spatial legs
contribute $E^2(\sin^2\theta\cos^2\varphi + \sin^2\theta\sin^2\varphi +
\cos^2\theta) = E^2$, giving $-E^2 + E^2 = 0$.
\implnote{The camera module constructs $k^\mu$ from a pixel grid and the
observer's tetrad, then hands the resulting null vector to the geodesic
integrator.}

The chapter introduction promised that this parametrisation of null
vectors would reappear verbatim in the camera model, and
\cref{eq:null-tetrad} is that parametrisation.

In curved spacetime, the tetrad vectors $e^\mu{}_{(a)}$ become functions of
position---they satisfy
$\metric\,e^\mu{}_{(a)}\,e^\nu{}_{(b)} = \eta_{ab}$ rather than the Minkowski
version---but \cref{eq:null-tetrad} itself is unchanged.  The flat-spacetime
machinery of this chapter is the complete template for relativistic ray tracing;
what remains is to let the metric vary.
\xrefnote{The geodesic equation, derived in
\cref{ch:differential-geometry}, governs how $k^\mu$ evolves as the ray
propagates through curved spacetime.}


\paragraph{Key References.}
The index-notation conventions follow \cite{MTW1973}; the tetrad formalism
draws on \cite{Wald1984}. The camera-tetrad construction is described in
\cite{James2015}.
